\documentclass[10pt,a4paper]{amsart}
\usepackage[margin=19mm]{geometry}
\usepackage{amsmath,amssymb,mathtools}
\usepackage[scr=rsfs,cal=euler]{mathalfa}
\usepackage{xfrac}
\usepackage[unicode,hidelinks,bookmarks=false]{hyperref}
\newcommand{\ie}{i.e.\ }
\newcommand{\cf}{cf.\ }
\newcommand{\resp}{respectively\ }
\newcommand{\Subsection}{\subsection}
\providecommand{\nobreakdash}{\nobreak\hbox{-}}
\setlength{\emergencystretch}{2em}
\multlinegap0pt
\usepackage{amssymb}
\usepackage{mathtools}
\usepackage{xcolor}
\usepackage[normalem]{ulem}
\usepackage{enumerate}
\newtheorem{prop}{Proposition}
\newtheorem{lem}{Lemma}
\def\crochet#1{\langle #1 \rangle}
\newcommand{\dF}{\ensuremath{\mathbb{F}}}
\title{Central limit theorems for martingales-II: convergence in the weak dual topology, with corrections}
\author{Bruno R\'{e}millard, Jean Vaillancourt}
\date{Source: arXiv:2304.04887v3 (CC BY 4.0)}
\begin{document}
\maketitle

\noindent\emph{Source and licence.} Central limit theorems for martingales-II: convergence in the weak dual topology, with corrections. Version arXiv:2304.04887v3 (CC BY 4.0), distributed by arXiv under the Creative Commons Attribution 4.0 International licence, http://creativecommons.org/licenses/by/4.0/, as recorded in the arXiv licence metadata for this exact version and shown on the abstract page https://arxiv.org/abs/2304.04887v3. Full licence notice: \texttt{licenses/arxiv-2304.04887-CC-BY-4.0.txt}.\par\smallskip

\noindent\textit{Editorial scope.} Counted unit: the article's prop prop:Cto0 (source0.tex lines 862--866). The article's complete original proof of it is delivered with it (source0.tex lines 868--886, one \texttt{proof} environment of 19 lines). No mathematics is written, repaired or re-derived: the statement and the proof are the article's own lines, in the article's own notation, and no retained line is substituted or re-flowed.  Retained uncounted as the source-local context the proof needs: lines 852--856; lines 936--956.  Nothing was omitted from any retained span, so the package carries no omission record.  No retained line contains a citation, so the wrapper carries no bibliography and no bibliography entry.  No retained line contains a figure environment, an image-inclusion command or drawing code, so no image is delivered and the package declares no asset dependency.  Editorial formatting only, in addition: an AMS \texttt{amsart} preamble that restates the article's own declarations and macros used by the retained lines, namely the environments \texttt{prop}, \texttt{lem} on separate counters, together with the standard packages those lines need.  The retained environments print in this document as Proposition 1, Lemma 1 in order of appearance, whereas the article prints the same items with the numbering of its own full document.  The complete native source of the article as distributed in the arXiv e-print for this exact version is bundled unchanged as \texttt{sources/139941/source0.tex}.
\begin{prop}\label{prop:Cincreasing}
Let $D$-valued non-decreasing processes $A_n$ and some continuous process $A$
be such that $A_n \stackrel{f.d.d.}{\rightsquigarrow} A$. Then $A_n$ is $\mathcal{C}$-tight
and $A_n \stackrel{\mathcal{C}}{\rightsquigarrow} A$.
\end{prop}

\begin{lem}[Lenglart's inequality]\label{lem:lenglart}
Let $X$ be an $\dF$-adapted $D$-valued process. Suppose that $Y$ is optional, non-decreasing,
and that, for any bounded stopping time $\tau$, $E|X(\tau)|\leq E\{Y(\tau)\}$.
Then for any stopping time $\tau$ and all $\varepsilon,\eta >0$,
\begin{itemize}
\item[{a)}]
if $Y$ is predictable,
\begin{equation}\label{eng1}
P(\sup_{s\leq \tau}|X(s)|\geq\varepsilon) \leq \frac{\eta}{\varepsilon} +
P(Y(\tau)\geq\eta).
\end{equation}

\item[b)]
if $Y$ is adapted,
\begin{equation}\label{eng2}
P(\sup_{s\leq \tau}|X(s)|\geq\varepsilon) \leq \frac{1}{\varepsilon}
\left[\eta+ E\left\{J_\tau(Y) \right\}\right] +
P(Y(\tau)\geq\eta).
\end{equation}
\end{itemize}
\end{lem}

\begin{prop}\label{prop:Cto0}
$M_n \stackrel{\mathcal{C}}{\rightsquigarrow} 0$,
$[M_n] \stackrel{\mathcal{C}}{\rightsquigarrow} 0$ and
$\crochet{ M_n }  \stackrel{\mathcal{C}}{\rightsquigarrow} 0$ are equivalent.
\end{prop}

\begin{proof} The operator on $D$ defined by mapping $x\in D$ to $t\mapsto\sup_{0\le s\le t}x(s)$,
is a continuous operator in the $\mathcal{C}$-topology. The three statements
$M_n \stackrel{\mathcal{C}}{\rightsquigarrow} 0$, $M_n^2 \stackrel{\mathcal{C}}{\rightsquigarrow} 0$ and
$\sup_{0\le s\le\{\cdot\}}M_n^2(s) \stackrel{\mathcal{C}}{\rightsquigarrow} 0$ are therefore equivalent.
By Proposition \ref{prop:Cincreasing}, it suffices to prove the proposition with each of
$M_n \stackrel{\mathcal{C}}{\rightsquigarrow} 0$,
 $[M_n] \stackrel{\mathcal{C}}{\rightsquigarrow} 0$ and
$\crochet{ M_n } \stackrel{\mathcal{C}}{\rightsquigarrow} 0$,
respectively replaced by $\sup_{0\le s\le t}M_n^2(s) \stackrel{Law}{\rightsquigarrow}0$,
$[M_n]_t\stackrel{Law}{\rightsquigarrow}0$ and
$\crochet{ M_n }_t\stackrel{Law}{\rightsquigarrow}0$, for all $t\ge0$ in all three cases.
Doob's inequality yields
$$
E\{[M_n]_t\} \le E\left\{\sup_{0\le s\le t}M_n^2(s)\right\}\le 4E\{[M_n]_t\}.
$$
Applying Lemma \ref{lem:lenglart} twice, interchanging the roles of $X$ and $Y$, yields the first equivalence for $M_n$ and $[M_n]$.
Equality $E\{\crochet{ M_n }_t\} = E\{[M_n]_t\}$ and another double application of Lemma \ref{lem:lenglart} yields the last one,
this time with $\crochet{ M_n }$ and $[M_n]$.
\end{proof}

\end{document}
